\documentclass{ximera}

\title{Module 4: Vertex Form Applications}

\author{Miller, Sims}

\begin{document}
\begin{abstract}
    Investigate how vertex form $y=a(x-h)2 +k$ models real-world maximum and minimum situations using projectile motion and optimization contexts. 
\end{abstract}

\subsection*{Scenario}
During a charity golf tournament at Prairie View A\&M University, a golfer is competing in a longest-drive contest. The golfer launches a golf ball from the tee, and the height of the ball above the ground can be modeled by the quadratic function. The golfer hits a golf ball, and its height above the ground is modeled by the quadratic function $$h(t)=-16(t-2)^2+80$$ where $h(t)$ represents the height of the golf ball in feet and $t$ represents the time in seconds after the ball is hit.

\subsection*{Student Task}
The tournament officials want to understand how high the golf ball travels and when it reaches its highest point. 

As a mathematical consultant for the tournament, your job is to analyze the trajectory of the golf ball using the vertex form of a quadratic function. 

\subsection*{Part A: Analyze the Function}
\begin{enumerate}
    \item Identify the vertex $(h,k)$ of the quadratic function. 

    \item Interpret the meaning of the vertex in the context of the problem. 

    \item Does the vertex represent a maximum or minimum value? Explain why. 

    \item What is the maximum height reached by the golf ball? 

    \item After how many seconds does the golf ball reach its maximum height? 
\end{enumerate}

\subsection*{Part B: Graph the Function}
\begin{enumerate}
    \item Graph the quadratic function using a graphing calculator, Desmos, or another graphing tool. 

    \item Clearly label: 
    \begin{itemize}
        \item the vertex
        \item the axis of symmetry 
        \item the y-intercept 
        \item the direction the parabola opens 
    \end{itemize}
    
    \item Include the graph in your project submission. 
\end{enumerate}

\subsection*{Part C: Mathematical Interpretation}
Answer the following questions in complete sentences. 
\begin{enumerate}
    \item What does the value $h=2$ tell us about the golf ball's flight?

    \item What does the value $k=80$ tell us about the golf ball's flight?

    \item Why is the leading coefficient $-16$ negative?

    \item Explain how vertex form helps us quickly identify important information about the golf ball's path. 
\end{enumerate}

\subsection*{Part D: Written Analysis}
Write a paragraph (5–8 sentences) addressing the following: 

Describe how the vertex of the quadratic function helps the golfer understand the trajectory of the golf ball. Explain the significance of the maximum height and the time at which it occurs. Discuss why vertex form is useful when modeling projectile motion. 


\subsection*{Part E: Presentation Component (OPTIONAL)}
Create one PowerPoint slide, Canva slide, or poster that includes: 
\begin{itemize}
    \item The quadratic equation 

    \item The graph 

    \item The vertex 

    \item A labeled diagram 

    \item A brief explanation of the golf ball's flight 
\end{itemize}

\subsection*{Extension Challenge}
A professional golfer wants the ball to reach a maximum height of 120 feet after 3 seconds. 

Write a new quadratic function in vertex form that models this situation and explain how you determined your answer. 
\end{document}